% Part: lambda-calculus
% Chapter: introduction
% Section: beta

\documentclass[../../../include/open-logic-section]{subfiles}

\begin{document}

\olfileid[pt]{lam}{int}{bet}
\olsection{Redução~$\beta$}

Quando vemos $(\lambd[m][(\lambd[y][y]) m])$, é natural conjecturar que ele
tem alguma relação com $\lambd[m][m]$, a saber, que o segundo termo deve ser o
resultado de «simplificar» o primeiro. A noção de \emph{redução~$\beta$}
captura formalmente essa intuição.

\begin{defn}[Contração~$\beta$, $\bredone$] \ollabel{defn:betacontr}
  A \emph{contração~$\beta$} ($\bredone$) é a menor relação compatível sobre
  termos que satisfaz a seguinte condição:
  \[
    (\lambd[x][N])Q \bredone \Subst{N}{Q}{x} 
  \]
  Dizemos que $P$ é \emph{contraído por~$\beta$} a $Q$ se $P \bredone Q$. Um
  termo da forma $(\lambd[x][N])Q$ é chamado de \emph{redex}.
\end{defn}

\begin{prob} \ollabel{prob:def}
  Explicite as definições indutivas equivalentes de contração~$\beta$, como
  fizemos para a mudança de variável ligada em
  \olref[lam][syn][alp]{defn:aconvone}.
\end{prob}
  
\begin{defn}[Redução~$\beta$, $\bred$] \ollabel{defn:betared}
  A \emph{redução~$\beta$} ($\bred$) é a menor relação reflexiva e transitiva
  sobre termos que contém~$\bredone$. Dizemos que $P$ é reduzido por~$\beta$
  a $Q$ se $P \bred Q$.
\end{defn}

Escreveremos $\redone$ em vez de $\bredone$, e $\red$ em vez de $\bred$,
quando o contexto for claro.

Em termos informais, $M \bred N$ se e somente se $M$ pode ser transformado em
$N$ por zero ou vários passos de contração~$\beta$.

\begin{defn}[Normal por~$\beta$]
Um termo que não pode mais ser contraído por~$\beta$ é dito
\emph{normal por~$\beta$}.
\end{defn}

Se $M \bred N$ e $N$ é normal por~$\beta$, dizemos que $N$ é uma
\emph{forma normal} de~$M$. Pode-se perguntar se a forma normal de um termo é
única, e a resposta é afirmativa, como veremos mais adiante.

Consideremos alguns exemplos.
\begin{enumerate}
\item Temos
\begin{align*}
(\lambd[x][xxy]) \lambd[z][z] & \redone (\lambd[z][z])(\lambd[z][z]) y \\
& \redone (\lambd[z][z]) y \\
& \redone y
\end{align*}
\item «Simplificar» um termo pode, na verdade, torná-lo mais complexo:
\begin{align*}
(\lambd[x][xxy])(\lambd[x][xxy]) & \redone (\lambd[x][xxy])(\lambd[x][xxy])y \\
& \redone (\lambd[x][xxy])(\lambd[x][xxy])yy \\
& \redone \dots
\end{align*}
\item Também pode deixar um termo inalterado:
\[
(\lambd[x][xx])(\lambd[x][xx]) \redone (\lambd[x][xx])(\lambd[x][xx])
\]
\item Além disso, alguns termos podem ser reduzidos de mais de uma maneira;
  por exemplo,
\[
(\lambd[x][(\lambd[y][yx]) z]) v \redone (\lambd[y][yv]) z
\]
pela contração da aplicação mais externa; e
\[
(\lambd[x][(\lambd[y][yx]) z]) v \redone (\lambd[x][zx]) v
\]
pela contração da aplicação mais interna. Observe, contudo, que, nesse caso,
ambos os termos se reduzem posteriormente ao mesmo termo, $zv$.
\end{enumerate}

O resultado final do último exemplo não é uma coincidência, mas ilustra uma
propriedade profunda e importante do cálculo lambda, conhecida como
propriedade de Church--Rosser.

\begin{digress}
  Em geral, há mais de uma maneira de reduzir um termo por~$\beta$; por isso,
  foram inventadas muitas estratégias de redução, entre as quais a mais comum
  é a \emph{estratégia natural}. A estratégia natural sempre contrai o redex
  \emph{mais à esquerda}, sendo a posição de um redex definida como seu ponto
  inicial no termo. A estratégia natural tem a propriedade útil de que um
  termo pode ser reduzido a uma forma normal por alguma estratégia se e
  somente se pode ser reduzido à forma normal mediante a estratégia natural.
  No que segue, usaremos a estratégia natural, salvo indicação em contrário.
\end{digress}

\begin{defn}[Equivalência~$\beta$, $\equal$]
  A \emph{equivalência~$\beta$} ($\equal$) é a relação definida
  indutivamente da seguinte maneira:
  \begin{enumerate}
  \item $M \equal M$.
  \item Se $M \equal N$, então $N \equal M$.
  \item Se $M \equal N$, $N \equal O$, então $M \equal O$.
  \item Se $M \equal N$, então $PM \equal PN$.
  \item Se $M \equal N$, então $MQ \equal NQ$.
  \item Se $M \equal N$, então $\lambd[x][M] \equal \lambd[x][N]$.
  \item $(\lambd[x][N])Q \equal \Subst{N}{Q}{x}$.
  \end{enumerate}
\end{defn}

As três primeiras regras tornam a relação uma relação de equivalência; as
três seguintes a tornam compatível; a última garante que ela contém a
contração~$\beta$.

Em termos informais, dois termos são $\beta$-equivalentes se e somente se um
deles pode ser transformado no outro por zero ou mais passos de
contração~$\beta$, ou pela «inversa» da contração~$\beta$. A inversa da
contração~$\beta$ é definida de modo que $M$ se contrai inversamente por
$\beta$ a $N$ se e somente se $N$ se contrai por~$\beta$ a $M$.

Além das regras acima, ampliaremos a relação com outras regras e denotaremos a
relação de equivalência ampliada por $\equal[X]$, onde $X$ é a regra de
extensão.

\end{document}
